\section{Results}

Sections~4.1--4.4 report canonical results; Section~4.5 adds post-hoc interpretation.

\subsection{Behaviour of the learned components}
\label{sec:components}

The CAE generalised to independent in-control data, with a mean squared reconstruction error per
window element of 0.66 on the D1 validation windows and 0.67 on D3. The OC-SVM met its
$\nu$-property in sample (3.0\% of D2 windows outside the boundary; 4.0\% support vectors), but
6.2\% of D3 windows fell outside, so its boundary was somewhat tight out of sample, which the
D3-calibrated limit absorbs. Figure~\ref{fig:ic} shows that the components were weakly correlated
($r=-0.15$), that $s_t$ was right-skewed (skewness 0.55) while $\ell_t$ was nearly symmetric, and
that both were strongly serially dependent: their lag-one autocorrelation of 0.96 decayed almost
linearly to near zero at lag 32, where windows stop sharing observations, whereas the raw variables'
lag-one autocorrelations of 0.45--0.60 fell below 0.07 by lag 10. This overlap-induced dependence is
why the limits were calibrated empirically.

\begin{figure}[!t]
\centering
\includegraphics[width=\linewidth]{diagnostic_ic_monitoring_vector.pdf}
\caption[In-control behaviour of the hybrid monitoring vector]{In-control monitoring vectors on
D3 ($300\times2969$ windows, recomputed from the frozen models; they reproduce the stored D3 moments
exactly). (a) Joint distribution of $(s_t,\ell_t)$ (grey shading: log count; cross: $\hat\mu$) with
marginal histograms. (b) Autocorrelation of $s_t$, $\ell_t$ and the four raw variables, pooled over
trajectories.}
\label{fig:ic}
\end{figure}

\subsection{In-control validation}

Table~\ref{tab:ic} sets each chart's D3 calibration value beside its independent D4 performance; no
D4 run length was censored (one D3 trajectory was censored when calibrating $H(0.20)$). The realised
$\ARL_0$ departed from 370 in opposite directions: the hybrid configurations ran shorter
(315.4--336.5, i.e.\ 2.3--4.0 D4 standard errors below target), the ablations stayed close (351.4 and
363.4) and $M(0.05)$ ran longer (403.5, 1.9 standard errors above); on the same D4 trajectories,
$M(0.05)$'s run lengths exceeded those of $H(0.05)$ by 67.0 on average (paired SE 22.2). Deviations
of this size can arise from sampling variability in both the finite calibration bank and D4, but
their direction matters: $H(0.05)$ was the more liberal chart in control, which, if anything,
favours it out of control.

\begin{table}[!t]
\centering
\caption[Calibrated limits and in-control validation]{Calibrated limit, D3 calibration $\ARL_0$ and
independent D4 in-control performance (500 trajectories; SE in parentheses) of the six charts.}
\label{tab:ic}
\begin{tabular}{llrrrrr}
\toprule
Chart & Input & $h^*$ & D3 $\ARL_0$ & D4 $\ARL_0$ (SE) & D4 SDRL & D4 MRL \\
\midrule
\tabinput{tables/ic_validation.tex}
\bottomrule
\end{tabular}
\end{table}

\subsection{Primary comparison: \texorpdfstring{$H(0.05)$ versus $M(0.05)$}{H(0.05) versus M(0.05)}}

Figure~\ref{fig:ooc}(a)--(b) covers the complete shift grid; no D5 run length was censored. At
$\delta=0.2$ the charts were indistinguishable ($\ARL_1$ 282.4 and 285.6; $\bar d=-3.2$, paired SE
15.9). At every $\delta\ge0.4$ the benchmark signalled first on average, by more than ten paired
standard errors (Appendix~\ref{app:results}). The absolute gap peaked at $\delta=0.4$--$0.6$
($\bar d\approx98$--106 opportunities) and then narrowed, while the relative gap widened: $M(0.05)$
signalled at its first opportunity in 41\% of trajectories at $\delta=1.0$ and in essentially all
from $\delta=2.2$, whereas $H(0.05)$ never did so and levelled off near an $\ARL_1$ of 12.5. Counted
in post-change observations ($j+31$), the mean delays were 110 against 43 at $\delta=1.0$ and 44
against 32 at $\delta=3.0$, a ratio of 1.4--2.5 for $\delta\ge0.4$, compared with up to 21 in
monitoring opportunities. Under this sustained directional mean shift, the classical benchmark
therefore detected faster at every shift from 0.4 upwards, despite being the more conservative chart
in control.

\begin{figure}[!t]
\centering
\includegraphics[width=\linewidth]{results_ooc_arl_panels.pdf}
\caption[Out-of-control performance across the shift grid]{Out-of-control results on D5 (500
trajectories per $\delta$; log scale in (a), (c), (d)). (a) $\ARL_1$ of $H(0.05)$ and $M(0.05)$;
(b) paired mean difference $\bar d(\delta)$ of Equation~(\ref{eq:paired}); (c) $H(0.05)$ and its
ablations; (d) the three hybrid smoothing constants, with $M(0.05)$ for reference.}
\label{fig:ooc}
\end{figure}

\subsection{Ablation and smoothing sensitivity}

Figure~\ref{fig:ooc}(c) shows that $H(0.05)$ had a lower $\ARL_1$ than both ablations at all 15
shifts. Paired over the same trajectories (Appendix~\ref{app:results}), its advantage over $R(0.05)$
was clear throughout (8.6--69.0 opportunities, more than three paired SE), whereas that over
$S(0.05)$ was within sampling error at $\delta\le0.4$ ($-9.7$, SE 14.5; $-6.4$, SE 11.3) and exceeded
two paired SE from $\delta=0.6$; the score was the stronger single component up to $\delta=1.4$ and
the reconstruction error beyond it. Figure~\ref{fig:ooc}(d) shows that a larger $\lambda$ changed
little at small shifts and helped at larger ones, as expected for EWMA-type charts
\citep{lowry_multivariate_1992}: $H(0.20)$ had the lowest hybrid $\ARL_1$ at every $\delta\ge0.4$, by
2.6--11.0 opportunities relative to $H(0.05)$. Its D4 $\ARL_0$ was also lower (paired difference
$-21.1$, SE 11.0), so part of this gain may reflect a more liberal realised calibration, and no
smoothing constant in the pre-specified range brought the hybrid close to $M(0.05)$.

\subsection{Interpretation}

Two features of the design explain these results. First, the fault is a sustained mean shift in a
fixed direction, the alternative that a linear, accumulating MEWMA on $x_t$ is built to detect
\citep{lowry_multivariate_1992}, whereas the hybrid accumulates two distance-type summaries of each
window, which respond weakly to small location shifts. Second, with $\tau=1$ the benchmark's first
report already contains 32 accumulated post-change updates, whereas the hybrid's first statistic is
the Mahalanobis distance of a single window, compared with a limit set by the strongly persistent
steady state. Post-hoc diagnostics, which are not part of the pre-specified experiment, support both
explanations: the component means moved by only 0.10--0.21 in-control standard deviations at
$\delta=0.4$ (Appendix~\ref{app:diagnostics}), and a separate diagnostic audit estimated that $y_t$
retains 2--36\% of the raw data's mean-shift information per update, that startup mechanics account
for most of the gap in monitoring opportunities at $\delta\ge1.8$, and that with the change after
100 in-control reports $H(0.05)$ still needed 2.0--2.5 times as many post-change observations as
$M(0.05)$ for $\delta\ge0.6$. These diagnostics explain the canonical result; they do not change it.
